\section{The Reach Readout}
\label{sec:method}

\begin{figure}[t]
\centering
\includegraphics[width=\linewidth]{fig_hero.png}
\caption{\textbf{The readout and the cross-regime inversion.} \emph{Left:} we displace the
object by $\dvec$ and measure the reach endpoint along $\uu$ (object $\uu{=}0$, canonical
$\uu{=}1$). A grounded reach lands near the object ($\uu{\to}0$); an anchored reach lands
near the canonical location ($\uu{\to}1$) yet can still succeed inside a wide grasp basin.
\emph{Right (data):} the \emph{same} counterfactual-fine-tuning intervention raises
displaced success in both suites, but spatial successes stay off-object
($\usucc\,1.06{\to}0.51$) while object successes become object-near
($\usucc\,0.67{\to}0.21$). Success rate cannot distinguish the two.}
\label{fig:hero}
\end{figure}

\paragraph{Setup.}
We study SmolVLA~\citep{smolvla2025} fine-tuned to competence on the LIBERO-spatial and
LIBERO-object suites~\citep{liu2023libero} (10 tasks each), reaching $74\%$ and $68\%$
clean success respectively. All rollouts run in simulation with fresh observations at
each planning step.

\paragraph{The intervention.}
At episode start we displace the task-relevant object by a planar vector $\dvec$ (a grid
of four magnitudes $\{25,50,75,100\}$\,mm $\times$ four directions), leaving the rest of
the scene intact, then roll out the policy under its normal chunked execution. The
canonical location $\Cpt$ is the object's training position; the displaced location is
$\Tpt = \Cpt + \dvec$.

\paragraph{The readout.}
Let $\Ept$ be a reach endpoint and $\evec = \Ept - \Tpt$ its error relative to the
displaced object. In the displacement-normalized frame we define
\begin{equation}
\uu = \frac{\langle \evec,\, -\hat{\dvec}\rangle}{\|\dvec\|},
\qquad
\vv = \frac{\langle \evec,\, \hat{\dvec}^{\perp}\rangle}{\|\dvec\|},
\qquad
\rrho = \frac{\|\evec\|}{\|\dvec\|},
\end{equation}
so $\uu = 0$ when the endpoint is at the object $\Tpt$ and $\uu = 1$ when it is at the
canonical location $\Cpt$; $\vv$ is the orthogonal component and $\rrho$ the normalized
error magnitude. We extract three endpoints per rollout: \emph{pre-grasp} (end-effector at
the first gripper-close command), \emph{contact} (first object motion), and \emph{lift}
(first object $z$-rise beyond a threshold). The pre-grasp endpoint captures the
\emph{planned} reach; contact and lift confirm whether the object was physically engaged.

\paragraph{Success conditioning.}
The defining move is to split $\uu,\vv,\rrho$ by rollout success. The question ``do
\emph{successful} spatial-generalization rollouts ground to the object?'' is exactly
$\usucc$: if displaced successes have $\usucc \approx 1$, the policy succeeds while reaching
the canonical location---basin tolerance, not grounding. We summarize a condition by the
success-conditioned gap $\ugap = \mathrm{median}(\ufail) - \mathrm{median}(\usucc)$:
$\ugap > 0$ means successful reaches are pulled toward the object more than failing ones.

\paragraph{Statistics.}
Endpoints are aggregated with a task-cluster (pairs) bootstrap over the 10 LIBERO tasks;
we report $90\%$ intervals. For ``unchanged'' claims we use two-one-sided-tests
equivalence with pre-declared margins. Full statistical details, including a
competence-matched degradation-null control used in \Cref{sec:boundary}, are in the
appendix.

\paragraph{Reusability.}
The readout needs only object displacement at reset and endpoint logging; it is agnostic
to the policy architecture and the task suite, and we release it as a drop-in tool.
